% Build with: latexmk -pdf main.tex
\documentclass[11pt]{article}

\usepackage[T1]{fontenc}
\usepackage{lmodern}
\usepackage[margin=1in]{geometry}
\usepackage{microtype}
\usepackage{amsmath,amssymb,amsthm,mathtools}
\usepackage{booktabs}
\usepackage{enumitem}
\usepackage{xcolor}
\usepackage[colorlinks=true,linkcolor=blue!50!black,citecolor=blue!50!black,urlcolor=blue!50!black]{hyperref}

\theoremstyle{definition}
\newtheorem{definition}{Definition}[section]

\newcommand{\bits}[1]{\{0,1\}^{#1}}
\newcommand{\getsr}{\stackrel{\$}{\gets}}
\newcommand{\Th}{\mathsf{TH}}
\newcommand{\IncEnc}{\mathsf{IncEnc}}
\newcommand{\KeyGen}{\mathsf{KeyGen}}
\newcommand{\Sign}{\mathsf{Sign}}
\newcommand{\Verify}{\mathsf{Verify}}
\newcommand{\Chain}{\mathsf{Chain}}
\newcommand{\PRF}{\mathsf{PRF}}
\newcommand{\Hash}{\mathsf{H}}
\newcommand{\LE}{\mathsf{LE}}
\newcommand{\Trunc}{\mathsf{Trunc}}
\newcommand{\concat}{\mathbin\Vert}
\newcommand{\ep}{\mathit{ep}}
\newcommand{\sk}{\mathit{sk}}
\newcommand{\pk}{\mathit{pk}}
\newcommand{\rootnode}{\mathit{root}}

\emergencystretch=1.5em

\title{XMSS for Ethereum consensus, specification draft.}
\author{}
\date{}

\begin{document}
\maketitle

\begin{abstract}
As proposed in~\cite{DKKW25}, XMSS is a promising post-quantum replacement for BLS, the signature scheme used by Ethereum validators. Aggregating XMSS requires a post-quantum snark. This technical note specifies an optimized XMSS instance with three goals: NIST security category~1~\cite{NISTPQC,DKKW25}, a signature smaller than the 1280-byte IPv6 minimum link MTU~\cite{RFC8200}, and snark-friendly encoding, optimized for leanVM-b~\cite{leanVMb}.
\end{abstract}

\section{Definitions and notation}

\begin{definition}[Synchronized signature scheme]
A synchronized signature scheme with message space $\mathcal M$ and lifetime $L$ consists of randomized algorithms $\KeyGen$ and $\Sign$ and a deterministic algorithm $\Verify$:
\[
  \KeyGen\longrightarrow(\pk,\sk),\qquad
  \Sign(\sk,\ep,m)\longrightarrow\sigma,\qquad
  \Verify(\pk,\ep,m,\sigma)\longrightarrow\{0,1\},
\]
where $m\in\mathcal M$ is the message and $\ep\in\{0,\ldots,L-1\}$ is the epoch. Correctness requires
\[
  \Pr\!\left[\Verify(\pk,\ep,m,\sigma)=1:(\pk,\sk)\gets\KeyGen,\ \sigma\gets\Sign(\sk,\ep,m)\right]=1
\]
for every $m\in\mathcal M$ and $\ep\in\{0,\ldots,L-1\}$. For each key pair, $\Sign(\sk,\ep,\cdot)$ may be invoked at most once for any fixed epoch $\ep$.
\end{definition}

XMSS is a synchronized signature scheme in which epoch $\ep$ selects a Winternitz one-time key authenticated by leaf $\ep$ of a Merkle tree~\cite{RFC8391}. This specification sets $\mathcal M=\bits{256}$ and $L=2^{32}$.

Byte strings are concatenated with $\concat$. Bits and integer encodings are little endian. $\LE_r(a)$ is the unsigned $r$-bit encoding of $a$. All indices are zero based.

\begin{center}
\begin{tabular}{@{}lll@{}}
\toprule
Symbol & Value & Meaning\\
\midrule
$n$ & $128$ bits & hash value and Merkle node length\\
$|s|$ & $256$ bits & master secret seed length\\
$|P|$ & $128$ bits & public hash parameter length\\
$|m|$ & $256$ bits & message length\\
$|\rho|$ & $192$ bits & encoding randomness length\\
$w$ & $3$ & log-length of Winternitz chains\\
$T$ & $194$ & required encoding digit sum\\
$h$ & $32$ & Merkle tree height\\
$L$ & $2^{32}$ & lifetime\\
\bottomrule
\end{tabular}
\end{center}

Let $\Hash:\bits{*}\to\bits{256}$ be a cryptographic hash function, and let $\Trunc_{128}$ keep the first 16 output bytes. Define
\[
  \Th(P,\tau,M)=\Trunc_{128}\!\left(\Hash(\tau\concat P\concat M)\right).
\]

\begin{definition}[Tweak encoding]
For a one-byte type $t$ and unsigned 32-bit integers $p$ and $j$, define the 16-byte tweak
\[
  \tau(t,p,j)=\LE_8(t)\concat\LE_{32}(p)\concat\LE_{32}(j)\concat 0^{56}.
\]
\end{definition}

The four public hash domains are fixed as follows.

\begin{center}
\begin{tabular}{@{}clll@{}}
\toprule
$t$ & Use & Sub-position $p$ & Index $j$\\
\midrule
$0$ & chain edge $k$ of chain $i$ & $8i+k$ & epoch $\ep$\\
$1$ & Winternitz public-key leaf & $0$ & epoch $\ep$\\
$2$ & Merkle node & parent level $\ell$ & node index $j$\\
$3$ & message encoding & $0$ & epoch $\ep$\\
\bottomrule
\end{tabular}
\end{center}

For a 32-byte master seed $s$, define
\[
  \PRF_s(d,a,b)=\Trunc_{128}\!\left(\Hash\!\left(s\concat\LE_{32}(d)\concat\LE_{64}(a)\concat\LE_{64}(b)\right)\right).
\]
The three domains are
\[
\begin{aligned}
  z_{\ep,i} &= \PRF_s(1000,\ep,i), &&\text{Winternitz secret chain start},\\
  P &= \PRF_s(1001,0,0), &&\text{public hash parameter},\\
  F_{\ell,j} &= \PRF_s(1002,\ell,j), &&\text{out-of-range Merkle filler}.
\end{aligned}
\]

The target-sum code is~\cite{DKKW25}
\[
  \mathcal C=\left\{(x_0,\ldots,x_{v-1})\in\{0,\ldots,7\}^{v}:\sum_{i=0}^{v-1}x_i=T\right\}.
\]
Define $\IncEnc(P,m,\rho,\ep)$ by computing
\[
  D=\Th\!\left(P,\tau(3,0,\ep),m\concat\rho\concat 0^{64}\right).
\]
Interpret the two halves of $D$ as unsigned little-endian words $d_0,d_1$. For $q\in\{0,1\}$ and $0\leq r<21$, set
\[
  x_{21q+r}=\left\lfloor\frac{d_q}{2^{3r}}\right\rfloor\bmod 8.
\]
Return $x=(x_0,\ldots,x_{41})$ if bit $63$ of both $d_0$ and $d_1$ is zero and $\sum_i x_i=194$; otherwise return $\bot$.

\begin{definition}[Hash chain]
For a chain index $i$, epoch $\ep$, starting position $k$, step count $a$, and value $z\in\bits{128}$, define $\Chain_{i,\ep}(P;k,a,z)$ recursively by
\[
  \Chain_{i,\ep}(P;k,0,z)=z
\]
and, when $k+a\leq6$,
\[
  \Chain_{i,\ep}(P;k,a+1,z)=\Th\!\left(P,\tau(0,8i+k+a,\ep),\Chain_{i,\ep}(P;k,a,z)\right).
\]
The valid range is $0\leq k\leq7$ and $0\leq a\leq7-k$.
\end{definition}

The public chain tips and their Merkle leaf are
\[
  q_{\ep,i}=\Chain_{i,\ep}(P;0,7,z_{\ep,i}),
  \qquad
  Q_{\ep}=\Th\!\left(P,\tau(1,0,\ep),q_{\ep,0}\concat\cdots\concat q_{\ep,41}\right),
  \qquad 0\leq i<42.
\]

\section{Key generation}

On input a nonempty interval $R=[\ep_{\min},\ep_{\max}]\subseteq\{0,\ldots,L-1\}$, $\KeyGen$ samples $s\getsr\bits{256}$ and derives $P$, the $z_{\ep,i}$, and the fillers $F_{\ell,j}$ as above. Let
\[
  I_{\ell,j}=\{j2^\ell,\ldots,(j+1)2^\ell-1\}
\]
and define
\[
  X_{0,j}=\begin{cases}
    Q_j,&j\in R,\\
    F_{0,j},&j\notin R.
  \end{cases}
\]
For $1\leq\ell\leq32$, define
\[
  X_{\ell,j}=\begin{cases}
    \Th\!\left(P,\tau(2,\ell,j),X_{\ell-1,2j}\concat X_{\ell-1,2j+1}\right),&I_{\ell,j}\cap R\neq\varnothing,\\
    F_{\ell,j},&I_{\ell,j}\cap R=\varnothing.
  \end{cases}
\]
Set $\rootnode=X_{32,0}$ and return
\[
  \pk=(\rootnode,P),
  \qquad
  \sk=(s,P,\ep_{\min},\ep_{\max}).
\]
The public key is serialized as $\rootnode\concat P$.

\section{Signing}

$\Sign(\sk,\ep,m)$ rejects if $\ep\notin[\ep_{\min},\ep_{\max}]$ or epoch $\ep$ has already been used. It samples $\rho\getsr\bits{192}$ until $x=\IncEnc(P,m,\rho,\ep)\neq\bot$, then computes
\[
  y_i=\Chain_{i,\ep}(P;0,x_i,z_{\ep,i}),
  \qquad 0\leq i<42,
\]
and the authentication path
\[
  A_\ell=X_{\ell,\,\left\lfloor\ep/2^\ell\right\rfloor\mathbin{\mathsf{xor}}1},
  \qquad 0\leq\ell<32.
\]
It returns
\[
  \sigma=(\rho,(y_0,\ldots,y_{41}),(A_0,\ldots,A_{31})).
\]

\section{Verification}

$\Verify(\pk,\ep,m,\sigma)$ parses $\pk=(\rootnode,P)$ and $\sigma=(\rho,(y_i)_i,(A_\ell)_\ell)$, then:
\begin{enumerate}[leftmargin=2em]
\item Compute $x=\IncEnc(P,m,\rho,\ep)$ and reject if $x=\bot$.
\item For every $0\leq i<42$, compute $q'_i=\Chain_{i,\ep}(P;x_i,7-x_i,y_i)$.
\item Set $Z_0=\Th(P,\tau(1,0,\ep),q'_0\concat\cdots\concat q'_{41})$.
\item For $0\leq\ell<32$, let $j=\lfloor\ep/2^{\ell+1}\rfloor$. If bit $\ell$ of $\ep$ is zero, set $Z_{\ell+1}=\Th(P,\tau(2,\ell+1,j),Z_\ell\concat A_\ell)$; otherwise set $Z_{\ell+1}=\Th(P,\tau(2,\ell+1,j),A_\ell\concat Z_\ell)$.
\item Accept if and only if $Z_{32}=\rootnode$.
\end{enumerate}

\section{Representation and cost}

The signature is serialized as
\[
  y_0\concat\cdots\concat y_{41}\concat\rho\concat A_0\concat\cdots\concat A_{31}.
\]
The public key is 32 bytes. The signature is $672+24+512=1208$ bytes. Fixed-size arrays have no length prefixes.

Verification performs 134 evaluations of $\Hash$: one encoding hash, $\sum_{i=0}^{41}(7-x_i)=42\cdot7-194=100$ chain hashes, one leaf hash, and 32 Merkle hashes. Excluding encoding retries and tree construction, signing performs $\sum_i x_i=194$ chain hashes.

\section{Security}

\emph{To be added.}

\bibliographystyle{alphaurl}
\bibliography{refs}

\end{document}
